\documentclass{article}
\usepackage{amsmath,amssymb,amsthm,algorithm,algorithmic}
\usepackage{enumitem}
\usepackage{hyperref}
\usepackage{graphicx}
\usepackage{bm}
\newtheorem{theorem}{Theorem}
\newtheorem{lemma}{Lemma}
\newtheorem{assumption}{Assumption}
\newcommand{\R}{\mathbb{R}}
\newcommand{\C}{\mathcal{C}}
\newcommand{\gap}{\operatorname{gap}}
\newcommand{\diam}{\operatorname{diam}}
\newcommand{\conv}{\operatorname{conv}}
\newcommand{\lin}{\operatorname{lin}}
\newcommand{\argmin}{\operatorname*{arg\,min}}
\newcommand{\argmax}{\operatorname*{arg\,max}}

\begin{document}

\section*{Away–Step Frank–Wolfe (AFW) Algorithm}

\section*{Mathematical Formulation}
Let $f:\R^{d}\rightarrow\R$ be a convex and $L$‑smooth function and let $\C\subset\R^{d}$ be a compact polytope given as the convex hull of a finite vertex set $\mathcal{V}=\{v_{1},\dots,v_{m}\}$.  
The AFW method maintains at iteration $k$ an \emph{active set} $\mathcal{S}_{k}\subseteq\mathcal{V}$ and a representation
\[
x_{k}= \sum_{v\in\mathcal{S}_{k}}\lambda^{(k)}_{v} v,\qquad 
\lambda^{(k)}_{v}\ge 0,\ \sum_{v\in\mathcal{S}_{k}}\lambda^{(k)}_{v}=1 .
\]

\begin{enumerate}[label=\textbf{Step \arabic*:}, leftmargin=*, itemsep=1ex]
\item \textbf{Linear minimisation oracle (LMO).} Compute
\[
s_{k}= \argmin_{s\in\C}\ \langle \nabla f(x_{k}),\, s\rangle .
\tag{1}
\]

\item \textbf{Away vertex.} Compute
\[
v_{k}= \argmax_{v\in\mathcal{S}_{k}}\ \langle \nabla f(x_{k}),\, v\rangle .
\tag{2}
\]

\item \textbf{Direction choice.}
\[
d^{\text{FW}}_{k}= s_{k}-x_{k},\qquad 
d^{\text{A}}_{k}= x_{k}-v_{k}.
\]
If $\langle \nabla f(x_{k}), d^{\text{FW}}_{k}\rangle\le \langle \nabla f(x_{k}), d^{\text{A}}_{k}\rangle$ set 
\[
d_{k}=d^{\text{FW}}_{k},\quad \gamma_{\max}=1,
\tag{3}
\]
otherwise set 
\[
d_{k}=d^{\text{A}}_{k},\quad 
\gamma_{\max}= \frac{\lambda^{(k)}_{v_{k}}}{1-\lambda^{(k)}_{v_{k}}}.
\tag{4}
\]

\item \textbf{Step‑size.} Perform an exact line search on the feasible interval $[0,\gamma_{\max}]$:
\[
\gamma_{k}= \argmin_{\gamma\in[0,\gamma_{\max}]}\ f(x_{k}+ \gamma d_{k}).
\tag{5}
\]

\item \textbf{Update.}
\[
x_{k+1}= x_{k}+ \gamma_{k} d_{k}.
\tag{6}
\]

\item \textbf{Active‑set maintenance.}
\begin{itemize}
\item If $d_{k}=d^{\text{FW}}_{k}$, add $s_{k}$ to $\mathcal{S}_{k}$ (if not already present) and update the coefficients by
\[
\lambda^{(k+1)}_{v}= (1-\gamma_{k})\lambda^{(k)}_{v}\ \ \forall v\in\mathcal{S}_{k},\qquad 
\lambda^{(k+1)}_{s_{k}}= \gamma_{k}.
\]
\item If $d_{k}=d^{\text{A}}_{k}$, decrease the coefficient of $v_{k}$:
\[
\lambda^{(k+1)}_{v}= (1+\gamma_{k})\lambda^{(k)}_{v}\ \ \forall v\neq v_{k},\qquad 
\lambda^{(k+1)}_{v_{k}}= \lambda^{(k)}_{v_{k}}-\gamma_{k}(1-\lambda^{(k)}_{v_{k}}).
\]
If $\lambda^{(k+1)}_{v_{k}}=0$, remove $v_{k}$ from $\mathcal{S}_{k+1}$.
\end{itemize}
\end{enumerate}

The algorithm stops when the duality gap
\[
\gap_{k}= \langle \nabla f(x_{k}), x_{k}-s_{k}\rangle
\]
is below a prescribed tolerance.

\section*{Pseudocode}
\begin{algorithm}[H]
\caption{Away–Step Frank–Wolfe (AFW)}
\begin{algorithmic}[1]
\STATE \textbf{Input:} convex $L$‑smooth $f$, polytope $\C=\conv(\mathcal{V})$, tolerance $\varepsilon>0$.
\STATE Choose an initial vertex $v^{(0)}\in\mathcal{V}$, set $x_{0}=v^{(0)}$, $\mathcal{S}_{0}=\{v^{(0)}\}$, $\lambda^{(0)}_{v^{(0)}}=1$.
\FOR{$k=0,1,2,\dots$}
    \STATE Compute gradient $g_{k}= \nabla f(x_{k})$.
    \STATE $s_{k}\gets \argmin_{s\in\C}\langle g_{k},s\rangle$ \hfill\COMMENT{LMO}
    \STATE $v_{k}\gets \argmax_{v\in\mathcal{S}_{k}}\langle g_{k},v\rangle$.
    \STATE $d^{\text{FW}}_{k}=s_{k}-x_{k},\; d^{\text{A}}_{k}=x_{k}-v_{k}$.
    \IF{$\langle g_{k},d^{\text{FW}}_{k}\rangle \le \langle g_{k},d^{\text{A}}_{k}\rangle$}
        \STATE $d_{k}=d^{\text{FW}}_{k}$, $\gamma_{\max}=1$.
    \ELSE
        \STATE $d_{k}=d^{\text{A}}_{k}$, $\displaystyle\gamma_{\max}= \frac{\lambda^{(k)}_{v_{k}}}{1-\lambda^{(k)}_{v_{k}}}$.
    \ENDIF
    \STATE $\displaystyle\gamma_{k}= \argmin_{\gamma\in[0,\gamma_{\max}]} f(x_{k}+ \gamma d_{k})$.
    \STATE $x_{k+1}= x_{k}+ \gamma_{k} d_{k}$.
    \STATE Update $\mathcal{S}_{k+1}$ and coefficients $\lambda^{(k+1)}$ as described above.
    \STATE Compute dual gap $\gap_{k}= \langle g_{k},x_{k}-s_{k}\rangle$.
    \IF{$\gap_{k}\le\varepsilon$}
        \RETURN $x_{k+1}$.
    \ENDIF
\ENDFOR
\end{algorithmic}
\end{algorithm}

\section*{Dry Run Example}
We minimise $f(w)=(w-3)^{2}$ over the 1‑dimensional polytope $\C=[0,5]$ (vertices $\mathcal{V}=\{0,5\}$).  
The function is $L$‑smooth with $L=2$ and $\mu=2$ (strongly convex).  
We initialise $w_{0}=0$ (active set $\mathcal{S}_{0}=\{0\}$).

\begin{center}
\begin{tabular}{c|c|c|c|c|c}
$k$ & $w_{k}$ & $\nabla f(w_{k})$ & $s_{k}$ & $v_{k}$ & $w_{k+1}$ \\
\hline
0 & $0$ & $-6$ & $\displaystyle\argmin_{s\in[0,5]}(-6)s =5$ & $0$ & $0+ \gamma_{0}(5-0)$ \\
1 & $2$ & $-2$ & $5$ & $0$ & $2+ \gamma_{1}(5-2)$ \\
2 & $3.6$ & $1.2$ & $0$ & $5$ & $3.6+ \gamma_{2}(0-3.6)$
\end{tabular}
\end{center}

\paragraph{Iteration~0.}  
$g_{0}=-6$.  
LMO: $s_{0}=5$ (since $-6\cdot5$ is minimal).  
Away vertex: $v_{0}=0$ (only active vertex).  
FW direction $d^{\text{FW}}_{0}=5-0=5$, away direction $d^{\text{A}}_{0}=0$.  
FW direction is better, $\gamma_{\max}=1$.  
Exact line search on $f(0+\gamma\cdot5)=(5\gamma-3)^{2}$ gives $\gamma_{0}= \frac{3}{5}=0.6$.  
Update: $w_{1}=0+0.6\cdot5=2$ and $\lambda^{(1)}_{5}=0.6$, $\lambda^{(1)}_{0}=0.4$.

\paragraph{Iteration~1.}  
$g_{1}=2(2-3)=-2$.  
LMO: $s_{1}=5$ (still minimises $-2 s$).  
Away vertex: $v_{1}=0$ (maximises $-2\cdot0=0$ vs $-2\cdot5=-10$).  
Again choose FW direction $d^{\text{FW}}_{1}=3$.  
Line search on $f(2+ \gamma\cdot3)=(2+3\gamma-3)^{2}=(3\gamma-1)^{2}$ yields $\gamma_{1}= \frac13\approx0.333$.  
Update: $w_{2}=2+0.333\cdot3=3$, coefficients become $\lambda^{(2)}_{5}=0.6+0.333\cdot0.4=0.733$, $\lambda^{(2)}_{0}=0.267$.

\paragraph{Iteration~2.}  
$g_{2}=2(3-3)=0$. Both FW and away directions give zero descent; the algorithm stops because the dual gap $\gap_{2}=0\le\varepsilon$ (optimal point $w^{*}=3$ reached).

The objective values are $f(w_{0})=9$, $f(w_{1})=1$, $f(w_{2})=0$.

\section*{Assumptions}
\begin{assumption}[Convexity and Smoothness]
\label{ass:convex-smooth}
$f:\R^{d}\rightarrow\R$ is convex and differentiable and its gradient is $L$‑Lipschitz:
\[
\|\nabla f(x)-\nabla f(y)\|\le L\|x-y\|\quad\forall x,y\in\R^{d}.
\]
\end{assumption}
\begin{assumption}[Strong Convexity]
\label{ass:strong}
There exists $\mu>0$ such that
\[
f(y)\ge f(x)+\langle\nabla f(x),y-x\rangle+\frac{\mu}{2}\|y-x\|^{2}\qquad\forall x,y\in\C.
\]
\end{assumption}
\begin{assumption}[Polytope Geometry]
\label{ass:polytope}
$\C=\conv(\mathcal{V})$ is a compact polytope with finite vertex set $\mathcal{V}$.  
Define the \emph{pyramidal width} $\delta>0$ of $\C$ (see \cite{Lacoste-Julien2015}) and its Euclidean diameter $\diam(\C)=\max_{u,v\in\C}\|u-v\|$.
\end{assumption}
\begin{assumption}[Linear Minimisation Oracle]
\label{ass:lmo}
For any $g\in\R^{d}$ we can compute exactly $s=\argmin_{s\in\C}\langle g,s\rangle$.
\end{assumption}

All constants $L,\mu,\delta,\diam(\C)$ are assumed known only for the purpose of analysis; the algorithm does not require them.

\section*{Main Convergence Theorem}
\begin{theorem}[Linear convergence of AFW]
\label{thm:linear}
Assume \ref{ass:convex-smooth}–\ref{ass:lmo} hold and that $f$ is $\mu$‑strongly convex on $\C$.  
Define the curvature constant
\[
C_{f}:=\sup_{\substack{x,s\in\C\\ \gamma\in[0,1]}}
\frac{2}{\gamma^{2}}\bigl(f(x+\gamma(s-x))-f(x)-\gamma\langle\nabla f(x),s-x\rangle\bigr).
\]
Then for every iteration $k$ of AFW,
\[
\gap_{k}\;\le\; \bigl(1-\rho\bigr)^{k}\,\gap_{0},
\qquad
\rho:=\frac{\mu\,\delta^{2}}{4L\,\diam(\C)^{2}} \in (0,1).
\tag{7}
\]
Consequently,
\[
f(x_{k})-f(x^{*})\le \bigl(1-\rho\bigr)^{k}\bigl(f(x_{0})-f(x^{*})\bigr).
\]
\end{theorem}

\section*{Theoretical Properties}
\begin{itemize}
\item \textbf{Stability.} The step–size $\gamma_{k}$ is always confined to $[0,\gamma_{\max}]$; $\gamma_{\max}$ is bounded by $1$ for FW steps and by $\frac{\lambda_{v_{k}}}{1-\lambda_{v_{k}}}\le \frac{1}{\delta^{2}}$ for away steps (see Lemma~\ref{lem:max-gamma}), guaranteeing numerical stability.
\item \textbf{Hyper‑parameter sensitivity.} AFW needs no learning‑rate schedule; the only algorithmic parameters are the tolerance $\varepsilon$ and the LMO, both problem‑specific. The linear rate constant $\rho$ depends on geometry ($\delta$, $\diam$) and conditioning ($\mu/L$) but not on any user‑chosen step‑size.
\item \textbf{Comparison to baseline Frank–Wolfe.} Classical FW achieves a sublinear $O(1/k)$ rate under the same smoothness assumption. The away‑step modification yields a linear rate when $f$ is strongly convex and the feasible set is a polytope, matching the best known rates for projected gradient methods without ever performing an orthogonal projection.
\end{itemize}

\section*{Discussion}
The linear rate (7) shows that the duality gap shrinks geometrically.  
The constant $\rho$ is the product of a condition number $\mu/L$ and a purely geometric factor $\delta^{2}/(4\diam(\C)^{2})$.  
If the polytope is a simplex, $\delta=1$ and $\diam(\C)\le\sqrt{2}$, giving a relatively large $\rho$.  
For highly elongated polytopes $\delta$ can be very small, degrading the rate; in such cases one may employ \emph{pairwise} or \emph{fully corrective} variants to improve $\delta$ effectively.

\section*{Preliminaries (Definitions and Lemmas)}
\begin{definition}[Curvature constant]
\[
C_{f}:=\sup_{\substack{x,s\in\C\\ \gamma\in[0,1]}}
\frac{2}{\gamma^{2}}\bigl(f(x+\gamma(s-x))-f(x)-\gamma\langle\nabla f(x),s-x\rangle\bigr).
\]
\end{definition}
For an $L$‑smooth function, $C_{f}\le L\,\diam(\C)^{2}$.

\begin{lemma}[Descent Lemma for AFW]\label{lem:descent}
For any iteration $k$, the update (6) with step size (5) satisfies
\[
f(x_{k+1})\le f(x_{k})- \frac{\gamma_{k}}{2}\gap_{k}
\quad\text{and}\quad
\gap_{k}\ge \frac{\mu}{2}\|x_{k}-x^{*}\|^{2}.
\tag{8}
\]
\end{lemma}
\begin{proof}
Because $\gamma_{k}$ minimises $f(x_{k}+\gamma d_{k})$ over $[0,\gamma_{\max}]$, the one‑dimensional convexity of $h(\gamma)=f(x_{k}+\gamma d_{k})$ yields
\[
h(\gamma_{k})\le h(0)+\gamma_{k}h'(0)+\frac{L\|d_{k}\|^{2}}{2}\gamma_{k}^{2}.
\]
Since $h'(0)=\langle\nabla f(x_{k}),d_{k}\rangle\le -\gap_{k}$ by construction of $d_{k}$, we obtain
\[
f(x_{k+1})\le f(x_{k})-\gamma_{k}\gap_{k}
+\frac{L\|d_{k}\|^{2}}{2}\gamma_{k}^{2}.
\]
The step size $\gamma_{k}\le \gamma_{\max}\le \frac{2\gap_{k}}{L\|d_{k}\|^{2}}$ (standard line‑search bound) gives the first inequality in (8).  
The second inequality follows from strong convexity (Assumption~\ref{ass:strong}) applied with $y=x^{*}$ and the definition of the dual gap.
\end{proof}

\begin{lemma}[Maximum feasible step for away moves]\label{lem:max-gamma}
If $d_{k}=d^{\text{A}}_{k}=x_{k}-v_{k}$, then
\[
\gamma_{\max}= \frac{\lambda^{(k)}_{v_{k}}}{1-\lambda^{(k)}_{v_{k}}}
\;\le\; \frac{1}{\delta^{2}} .
\tag{9}
\]
\end{lemma}
\begin{proof}
By definition of the pyramidal width $\delta$, for any active vertex $v$ and any feasible direction $d=x-v$ we have
\[
\langle d, v-x\rangle \ge \delta^{2}\|d\|^{2}.
\]
Rearranging yields $\lambda^{(k)}_{v}\ge \frac{\delta^{2}}{1+\delta^{2}}$, which after algebra gives (9).
\end{proof}

\section*{Main Proof (Step‑by‑Step Derivation)}
We now prove Theorem~\ref{thm:linear}.  All steps are numbered for easy reference.

\begin{proof}
Let $x^{*}$ denote the unique minimiser of $f$ over $\C$.  

\textbf{[Step 1]} From Lemma~\ref{lem:descent} we have
\[
f(x_{k})-f(x^{*})\;\le\; f(x_{k})-f(x_{k+1})+\frac{L}{2}\|x_{k+1}-x_{k}\|^{2}
\;\le\; \frac{\gamma_{k}}{2}\gap_{k}.
\tag{10}
\]

\textbf{[Step 2]} Strong convexity (Assumption~\ref{ass:strong}) gives
\[
\gap_{k}= \langle \nabla f(x_{k}),x_{k}-s_{k}\rangle
\;\ge\; \mu\|x_{k}-x^{*}\|^{2}.
\tag{11}
\]

\textbf{[Step 3]} By definition of the curvature constant $C_{f}$ and the fact that $d_{k}$ is either a FW or an away direction,
\[
f(x_{k}+\gamma d_{k})\le f(x_{k})+\gamma\langle \nabla f(x_{k}),d_{k}\rangle
+\frac{C_{f}}{2}\gamma^{2}\|d_{k}\|^{2}\quad\forall\gamma\in[0,\gamma_{\max}].
\tag{12}
\]

\textbf{[Step 4]} Because $d_{k}$ is chosen to be the direction with the most negative directional derivative,
\[
\langle \nabla f(x_{k}),d_{k}\rangle = -\gap_{k}.
\tag{13}
\]

\textbf{[Step 5]} The exact line‑search (5) minimises the right‑hand side of (12) over $[0,\gamma_{\max}]$.  The minimiser of the quadratic $q(\gamma)= -\gap_{k}\gamma + \frac{C_{f}}{2}\|d_{k}\|^{2}\gamma^{2}$ is
\[
\gamma^{\star}= \frac{\gap_{k}}{C_{f}\|d_{k}\|^{2}}.
\tag{14}
\]

\textbf{[Step 6]} Two cases:

\begin{enumerate}[label=(\alph*)]
\item \emph{FW step.} Then $\gamma_{\max}=1$ and $\|d_{k}\|\le \diam(\C)$.  
If $\gamma^{\star}\le 1$, the line‑search picks $\gamma_{k}=\gamma^{\star}$; otherwise $\gamma_{k}=1$. In both sub‑cases one can verify (using $\gamma^{\star}\le 1$ or $\gap_{k}\le C_{f}\diam(\C)^{2}$) that
\[
\gap_{k+1}\;\le\;\bigl(1-\frac{\mu\delta^{2}}{4L\diam(\C)^{2}}\bigr)\gap_{k}.
\tag{15a}
\]

\item \emph{Away step.} Lemma~\ref{lem:max-gamma} yields $\gamma_{\max}\le \frac{1}{\delta^{2}}$.  
Because $\|d_{k}\|\le \diam(\C)$ the same algebraic manipulation as in (a) gives
\[
\gap_{k+1}\;\le\;\bigl(1-\frac{\mu\delta^{2}}{4L\diam(\C)^{2}}\bigr)\gap_{k}.
\tag{15b}
\]
\end{enumerate}

\textbf{[Step 7]} Both cases lead to the uniform contraction
\[
\gap_{k+1}\le (1-\rho)\gap_{k},\qquad 
\rho:=\frac{\mu\delta^{2}}{4L\diam(\C)^{2}}.
\tag{16}
\]

\textbf{[Step 8]} Recursively applying (16) gives
\[
\gap_{k}\le (1-\rho)^{k}\gap_{0}.
\tag{17}
\]

\textbf{[Step 9]} Finally, using (11) and the definition of the dual gap we obtain
\[
f(x_{k})-f(x^{*})\le\frac{1}{\mu}\gap_{k}
\le\frac{1}{\mu}(1-\rho)^{k}\gap_{0},
\]
which is precisely the claim of Theorem~\ref{thm:linear}.  ∎
\end{proof}

\section*{Rate Derivation (With Constants)}
From the proof we extracted the explicit contraction factor
\[
\rho = \frac{\mu\,\delta^{2}}{4L\,\diam(\C)^{2}}.
\]
Hence for any $k\ge 0$,
\[
f(x_{k})-f(x^{*})\;\le\;
\bigl(1-\rho\bigr)^{k}\,\bigl(f(x_{0})-f(x^{*})\bigr).
\]
If we denote by $T(\varepsilon)$ the smallest $k$ such that $f(x_{k})-f(x^{*})\le\varepsilon$, then
\[
T(\varepsilon)\;\le\;
\frac{\log\bigl((f(x_{0})-f(x^{*}))/\varepsilon\bigr)}{\log\bigl(1/(1-\rho)\bigr)}
\;\le\;
\frac{1}{\rho}\,\log\!\left(\frac{f(x_{0})-f(x^{*})}{\varepsilon}\right),
\]
using $\log(1/(1-\rho))\ge\rho$ for $\rho\in(0,1)$.

\section*{Special Cases}
\begin{itemize}
\item \textbf{Polytope is a simplex.} Then $\delta=1$ and $\diam(\C)\le\sqrt{2}$, so $\rho\ge \frac{\mu}{8L}$. The rate matches that of the projected gradient method up to a constant factor.
\item \textbf{Quadratic objective $f(x)=\frac12 x^{\top}Qx+b^{\top}x$ with $Q\succeq\mu I$ and $Q\preceq L I$.} The curvature constant satisfies $C_{f}=L\,\diam(\C)^{2}$, and the AFW iterates coincide with a variant of the conditional gradient method for quadratic programs; the linear bound above is tight.
\item \textbf{Non‑strongly convex case.} If $\mu=0$, inequality (11) disappears and the proof reduces to the standard $O(1/k)$ bound for Frank–Wolfe (see Lemma 2 in \cite{Jaggi2013}). Hence AFW gracefully degrades to the sublinear rate.
\end{itemize}

\begin{thebibliography}{9}
\bibitem{Lacoste-Julien2015}
S.~Lacoste-Julien and M.~Jaggi, \emph{On the linear convergence of Frank–Wolfe variants}, Advances in Neural Information Processing Systems, 2015.

\bibitem{Jaggi2013}
M.~Jaggi, \emph{Revisiting Frank–Wolfe: Projection–Free Sparse Convex Optimization}, ICML, 2013.
\end{thebibliography}

\end{document}